\chapter{Канонический аффинно-хиральный подъём и устранение двойного семейного фактора}

\section{Проверяемая щель}

Предыдущий гейт получил хиральный функционал на одном пакете
$H_L^8\to H_R^7$, но оставил открытой карту из фундаментального поля Тома VII.
Исходно было записано
\begin{equation}
 \Phi\in E_{\rm aff}\widehat\otimes\mathcal Y_{\rm phys},
 \qquad
 \mathcal Y_{\rm phys}=\mathcal E_\rho\otimes\Lambda_{\rm ch}.
 \label{eq:v7-original-double-family-carrier}
\end{equation}
Здесь
\begin{equation}
 \dim_{\mathbb C}E_{\rm aff}=12,
 \qquad
 \dim_{\mathbb C}\mathcal E_\rho=4,
 \qquad
 \dim_{\mathbb C}\Lambda_{\rm ch}=3.
\end{equation}
Но $E_{\rm aff}$ уже является динамическим аффинным семейным носителем, а
$\mathcal E_\rho$ --- касательным пространством после выбора конкретного
ранг-один семейного проектора $\rho$. Их тензорное произведение требует
проверки на двойной счёт и ковариантность.

\section{Неинвариантность фиксированного касательного пространства}

Для $\rho=|e_1\rangle\langle e_1|$ имеем
\begin{equation}
 \mathcal E_\rho
 =\rho M_3Q\oplus QM_3\rho,
 \qquad Q=I_3-\rho.
 \label{eq:v7-fixed-rho-tangent}
\end{equation}
Это четырёхмерное комплексное пространство является касательным в одной
точке ранговой орбиты, а не глобальным линейным представлением полного
$S_4$. Машинный аудит сопрягает его всеми двадцатью четырьмя матрицами
стандартного триплета. Восемьдесят из девяноста шести образов базисных
матриц выходят из $\mathcal E_\rho$, максимальный ортогональный остаток равен
\begin{equation}
 \frac{\sqrt3}{2}=0.8660254037\ldots,
\end{equation}
а линейная оболочка всей орбиты имеет размерность восемь.

\begin{theorem}[Запрет фиксированного $\mathcal E_\rho$ в нулевом поле]
Фиксированное касательное пространство $\mathcal E_\rho$ нельзя использовать
как глобальный $S_4$-ковариантный линейный множитель фундаментального поля,
если сам проектор $\rho$ должен возникать из этого поля после перехода.
Такое использование заранее выбирает точку будущей орбиты.
\end{theorem}

Есть и независимый линейный тест. Пространство
\begin{equation}
 E_{\rm aff}=\operatorname{Hom}(\mathbb C^4,\operatorname{im}P_3)
\end{equation}
имеет ровно одно $S_4$-инвариантное направление, порождённое канонической
коизометрией $V$. Поэтому любой инвариантный линейный функционал на
$E_{\rm aff}$ пропорционален
\begin{equation}
 \ell_V(X)=\Tr(V^\dagger X).
 \label{eq:v7-affine-invariant-functional}
\end{equation}
Свёртка исходного поля через $\ell_V$ имеет ядро размерности одиннадцать на
аффинном множителе и уничтожает $11\cdot12=132$ комплексных направлений.
Она не является полным динамическим подъёмом.

\section{Минимальная исправленная типизация}

Динамический аффинный модуль должен заменить локальное
$\mathcal E_\rho$, а не умножаться на него. Поэтому минимальное поле
кратностей имеет вид
\begin{equation}
 \Psi\in E_{\rm aff}\widehat\otimes\Lambda_{\rm ch},
 \qquad
 \dim_{\mathbb C}=12\cdot3=36.
 \label{eq:v7-corrected-affine-edge-field}
\end{equation}
В принятой конвенции, не пересчитывающей компоненты самого хиггсовского
дублета, это семьдесят два вещественных направления.

Тензорное тождество даёт каноническое включение
\begin{equation}
 E_{\rm aff}\otimes\operatorname{Hom}(H_L,H_R)
 \simeq
 \operatorname{Hom}(\mathbb C^4\otimes H_L,
                     \operatorname{im}P_3\otimes H_R).
 \label{eq:v7-canonical-tensor-hom-lift}
\end{equation}
Три физических ребра задают подмодуль правой части без выбора функционала
на $E_{\rm aff}$.

\section{Поднятый хиральный функционал}

Физический левый проектор и правый носитель равны
\begin{equation}
 P_L=P_3\otimes I_8,
 \qquad
 \rank P_L=24,
 \qquad
 H_R^{(3)}=\operatorname{im}P_3\otimes H_R,
 \qquad
 \dim H_R^{(3)}=21.
\end{equation}
Для поднятого нечётного оператора $Z$ положим
\begin{equation}
 \mathcal S_{\rm lift}(Z)
 =\frac1{45}\left\{
 \Tr_{32}(P_L-Z^\dagger Z)^2
 +\Tr_{21}(ZZ^\dagger-I_{21})^2
 \right\}.
 \label{eq:v7-lifted-hodge-parent}
\end{equation}
Поле удовлетворяет $Z=ZP_L$, поэтому восемь равномерных левых состояний
$P_1\otimes H_L$ являются опорным углом и не входят в физический ранг.

Если $\sigma_1,\ldots,\sigma_{21}$ --- сингулярные числа ограничения $Z$ на
$\operatorname{im}P_L$, то
\begin{equation}
 \mathcal S_{\rm lift}(Z)
 =\frac1{45}\left[
 3+2\sum_{j=1}^{21}(1-\sigma_j^2)^2
 \right]\ge\frac1{15}.
 \label{eq:v7-lifted-singular-value-action}
\end{equation}
В нуле
\begin{equation}
 \mathcal S_{\rm lift}(Z)
 =1-\frac4{45}\|Z\|_{\rm HS}^2+O(\|Z\|^4),
 \qquad
 \Hess_0\mathcal S_{\rm lift}=-\frac8{45}I
 \label{eq:v7-lifted-negative-hessian}
\end{equation}
на ортонормированном физическом касательном пространстве. В частности, все
семьдесят два вещественных направления исправленного модуля кратностей
отрицательны.

\section{Явный минимум и три ядерные линии}

Пусть $V:\mathbb C^4\to\operatorname{im}P_3$ --- каноническая аффинная
коизометрия, а $Y_\star:H_L^8\to H_R^7$ --- физическая коизометрия рёбер
$u,d,e$ предыдущего гейта. Тогда
\begin{equation}
 Z_\star=V\otimes Y_\star
 \label{eq:v7-factorized-three-generation-vacuum}
\end{equation}
принадлежит исправленному физическому модулю и удовлетворяет
\begin{equation}
 Z_\star Z_\star^\dagger=I_{21},
 \qquad
 Z_\star^\dagger Z_\star=P_3\otimes P_7,
 \qquad
 \rank Z_\star=21.
\end{equation}
Полное ядро в $\mathbb C^4\otimes H_L$ имеет размерность одиннадцать и
канонически раскладывается:
\begin{equation}
 \ker Z_\star
 =(\operatorname{im}P_1\otimes H_L)
 \oplus
 (\operatorname{im}P_3\otimes\ker Y_\star),
 \qquad
 11=8+3.
 \label{eq:v7-kernel-reference-physical-split}
\end{equation}
После ограничения на физический угол остаются ровно три комплексные ядерные
линии. Это по одной линии на каждое аффинное поколение, но ещё не вывод
нейтринных масс или смешивания.

\begin{theorem}[Канонический трёхпоколенный подъём]
Исправленный модуль $E_{\rm aff}\otimes\Lambda_{\rm ch}$ допускает
каноническое включение в трёхпоколенный хиральный оператор. Поднятый
Hodge-функционал имеет стационарный неустойчивый нуль, ограничен снизу
значением $1/15$ и достигает его на разрешённом операторе ранга $21$.
Физическое ядро минимума имеет комплексную размерность три.
\end{theorem}

\section{Что остаётся открытым}

Исправление закрывает типизацию, но не всё физическое замыкание.
Необходимо:
\begin{enumerate}
 \item построить полный Real-оператор и проверить junk/BRST--BV фактор;
 \item определить вакуумное многообразие внутри всего разрешённого, не
 обязательно факторизованного модуля;
 \item проверить, фиксирует ли функционал две относительные амплитуды
 рёбер $u,d,e$ и какие безмассовые орбитальные направления остаются;
 \item вывести пространственную жёсткость и динамику роста отрицательных мод.
\end{enumerate}

\begin{nogocriterion}[Запрет возврата к размерности 144]
Нельзя одновременно считать $E_{\rm aff}$ динамической заменой выбранного
семейного проектора и сохранять фиксированный касательный модуль
$\mathcal E_\rho$ как независимый линейный множитель. Возврат к
$12\cdot12=144$ комплексным компонентам требует нового глобального
расслоения касательных пространств и отдельной меры.
\end{nogocriterion}

\begin{protocol}
\begin{enumerate}
 \item исходный модуль размерности $144$: \textbf{закрыт как двойной
 семейный счёт};
 \item инвариантных направлений в $E_{\rm aff}$: \textbf{$1$};
 \item ядро инвариантной линейной свёртки: \textbf{$11$};
 \item орбитальная оболочка фиксированного $\mathcal E_\rho$:
 \textbf{размерность $8$, не $4$};
 \item исправленный комплексный модуль кратностей: \textbf{$36$};
 \item исправленный вещественный касательный размер: \textbf{$72$};
 \item поднятый физический ранг: \textbf{$21$};
 \item опорное ядро: \textbf{$8$};
 \item физическое ядро: \textbf{$3$};
 \item минимальная энергия: \textbf{$1/15$};
 \item следующий тест: \textbf{полный вакуумный гессиан исправленного
 Real-модуля и классификация его минимумов}.
\end{enumerate}
\end{protocol}

Машинный сертификат сохраняется в
\path{s2t_v7_affine_physical_module_canonical_lift_gate_results.json}.